\documentclass[10pt,a4paper]{amsart}
\usepackage{amsmath,amssymb,amsthm,graphicx}
% Retained minimal preamble for the Stacks Project chapter more-algebra.tex (More on Algebra).
% The chapter file pinned-original-source.tex carries only \input{preamble} on its first line, so no
% preamble text lives inside the chapter itself.  The Stacks Project uses one single shared preamble for
% every chapter; the snapshot stages it identically as sources/stacks-chow/preamble.tex and
% sources/stacks-etale-cohomology/preamble.tex.  The block below is copied mechanically, line for line,
% from sources/stacks-chow/preamble.tex lines 14-18 (the Xy-pic setup) and 154-217 (the native theorem
% declarations together with the mathematical macros).  Nothing was added, removed or redefined; the
% wrapper class and the amsmath/amssymb/amsthm/graphicx base packages are supplied by the assembler.
% Chapter: more-algebra.tex -- More on Algebra.
% For commutative diagrams we use Xy-pic
\usepackage[all]{xy}

% We use 2cell for 2-commutative diagrams.
\xyoption{2cell}

% Theorem environments.
%
\theoremstyle{plain}
\newtheorem{theorem}[subsection]{Theorem}
\newtheorem{proposition}[subsection]{Proposition}
\newtheorem{lemma}[subsection]{Lemma}

\theoremstyle{definition}
\newtheorem{definition}[subsection]{Definition}
\newtheorem{example}[subsection]{Example}
\newtheorem{exercise}[subsection]{Exercise}
\newtheorem{situation}[subsection]{Situation}

\theoremstyle{remark}
\newtheorem{remark}[subsection]{Remark}
\newtheorem{remarks}[subsection]{Remarks}

\numberwithin{equation}{subsection}

% Macros
%
\def\lim{\mathop{\mathrm{lim}}\nolimits}
\def\colim{\mathop{\mathrm{colim}}\nolimits}
\def\Spec{\mathop{\mathrm{Spec}}}
\def\Hom{\mathop{\mathrm{Hom}}\nolimits}
\def\Ext{\mathop{\mathrm{Ext}}\nolimits}
\def\SheafHom{\mathop{\mathcal{H}\!\mathit{om}}\nolimits}
\def\SheafExt{\mathop{\mathcal{E}\!\mathit{xt}}\nolimits}
\def\Sch{\mathit{Sch}}
\def\Mor{\mathop{\mathrm{Mor}}\nolimits}
\def\Ob{\mathop{\mathrm{Ob}}\nolimits}
\def\Sh{\mathop{\mathit{Sh}}\nolimits}
\def\NL{\mathop{N\!L}\nolimits}
\def\CH{\mathop{\mathrm{CH}}\nolimits}
\def\proetale{{pro\text{-}\acute{e}tale}}
\def\etale{{\acute{e}tale}}
\def\QCoh{\mathit{QCoh}}
\def\Ker{\mathop{\mathrm{Ker}}}
\def\Im{\mathop{\mathrm{Im}}}
\def\Coker{\mathop{\mathrm{Coker}}}
\def\Coim{\mathop{\mathrm{Coim}}}

% Boxtimes
%
\DeclareMathSymbol{\boxtimes}{\mathbin}{AMSa}{"02}

%
% Macros for moduli stacks/spaces
%
\def\QCohstack{\mathcal{QC}\!\mathit{oh}}
\def\Cohstack{\mathcal{C}\!\mathit{oh}}
\def\Spacesstack{\mathcal{S}\!\mathit{paces}}
\def\Quotfunctor{\mathrm{Quot}}
\def\Hilbfunctor{\mathrm{Hilb}}
\def\Curvesstack{\mathcal{C}\!\mathit{urves}}
\def\Polarizedstack{\mathcal{P}\!\mathit{olarized}}
\def\Complexesstack{\mathcal{C}\!\mathit{omplexes}}
% \Pic is the operator that assigns to X its picard group, usage \Pic(X)
% \Picardstack_{X/B} denotes the Picard stack of X over B
% \Picardfunctor_{X/B} denotes the Picard functor of X over B
\def\Pic{\mathop{\mathrm{Pic}}\nolimits}
\def\Picardstack{\mathcal{P}\!\mathit{ic}}
\def\Picardfunctor{\mathrm{Pic}}
\def\Deformationcategory{\mathcal{D}\!\mathit{ef}}
\usepackage{cleveref}
\crefname{originallemma}{original lemma}{original lemmas}
\Crefname{originallemma}{Original lemma}{Original lemmas}
\makeatletter
\newlabel{fields-lemma-Artin-Schreier}{{09DY}{source}{}{sourcepointer-1}{}}
\newlabel{fields-lemma-Artin-Schreier@cref}{{[originallemma][0][]09DY}{[1][0][]source}{}{}{}}
\makeatother
\crefname{originallemma}{original lemma}{original lemmas}
\Crefname{originallemma}{Original lemma}{Original lemmas}
\makeatletter
\newlabel{lemma-cohen}{{34747}{source}{}{sourcepointer-2}{}}
\newlabel{lemma-cohen@cref}{{[originallemma][0][]34747}{[1][0][]source}{}{}{}}
\makeatother
\crefname{originallemma}{original lemma}{original lemmas}
\Crefname{originallemma}{Original lemma}{Original lemmas}
\makeatletter
\newlabel{lemma-make-degree-q-extension}{{34641}{source}{}{sourcepointer-3}{}}
\newlabel{lemma-make-degree-q-extension@cref}{{[originallemma][0][]34641}{[1][0][]source}{}{}{}}
\makeatother
\crefname{originallemma}{original lemma}{original lemmas}
\Crefname{originallemma}{Original lemma}{Original lemmas}
\makeatletter
\newlabel{lemma-pre-characteristic-p-case}{{34765}{source}{}{sourcepointer-4}{}}
\newlabel{lemma-pre-characteristic-p-case@cref}{{[originallemma][0][]34765}{[1][0][]source}{}{}{}}
\makeatother
\crefname{originallemma}{original lemma}{original lemmas}
\Crefname{originallemma}{Original lemma}{Original lemmas}
\makeatletter
\newlabel{lemma-weakly-unramified-goes-up-along-totally-ramified}{{34490}{source}{}{sourcepointer-5}{}}
\newlabel{lemma-weakly-unramified-goes-up-along-totally-ramified@cref}{{[originallemma][0][]34490}{[1][0][]source}{}{}{}}
\makeatother
\title{In characteristic p, weak solvability for a degree p Galois step follows when the small residue field is the intersection of the p-power residue subfields}
\author{The Stacks Project Authors}
\date{stacks-project:more-algebra.tex:local-snapshot-20261012; GNU FDL 1.2 or later}
\begin{document}
\maketitle
\setcounter{subsection}{1}
\noindent\textit{Editorial note.} All delivered mathematical text is the author's own, copied from the pinned Stacks Project chapter file more-algebra.tex; no statement or proof step was rewritten, reordered or removed.\par\medskip
\noindent\textit{Editorial note.} The chapter file pinned-original-source.tex carries only \textbackslash{}input\{preamble\} on its first line, so no preamble text lives inside the chapter. The retained file sources/stacks-more-algebra/retained-preamble.tex is a mechanical, line-for-line copy of the project's single shared preamble as staged in the snapshot; nothing was added, removed or redefined and no package was added.\par\medskip
\noindent\textit{Editorial note.} The source URL is the source-level Stacks Project URL, as required when more than one card is produced from one source; the per-claim tag is carried by claim\_key.\par\medskip
\noindent\textit{Editorial note.} The author defines J by the condition lambda\_i not in kappa\_A while indexing it with i, writing 'J = \{j in \{n, ..., 1\} | lambda\_i not in kappa\_A\}' with a mismatched symbol; the definition is reproduced unchanged. The footnote about the case where B or A is complete, and the parenthetical '(Please observe that in this procedure we may get thrown back into the case that J = empty we treated above.)', are both inside the counted proof span and are delivered verbatim.\par\medskip
\noindent\textit{Allowed standard prerequisites (editorial).}\begin{itemize}
\item Artin-Schreier theory in characteristic p. invoked at proof line 34833 and used to parameterise M over L. Gives the element z with xi = z\textasciicircum{}p - z in L whose minimal n and later changes of choice drive the whole argument.
\item Cohen's structure theorem producing a coefficient map kappa\_B -> B. invoked at proof lines 34843-34848. Supplies sigma, in terms of which xi is expanded as a sum over powers of pi.
\item Weakly unramified extensions and their behaviour under totally ramified base change. assumption (1) at statement line 34822 and used at proof lines 34878-34880 and 34931-34933. Is what may be preserved when passing to K\_1 and L\_1, so that the weak-solution conclusion can be transported.
\item Construction of totally ramified degree q extensions with prescribed uniformizer. invoked at proof lines 34926-34930. Produces the auxiliary extension K\_1 of degree p\textasciicircum{}r used in the second half of the proof.
\end{itemize}\medskip
\ifdefined\hypertarget\hypertarget{sourcepointer-1}{}\fi
\noindent\textit{Source reference.} Original lemma 09DY: focused-seeds/AGP/sources/stacks-snapshot-20261011/tags.txt, source lines 11881--11881. Cross-chapter navigation pointer into the Stacks Project; the standard fact it supplies is applied inside the counted proof span and the argument resting on it stays there.\par\medskip
\ifdefined\hypertarget\hypertarget{sourcepointer-2}{}\fi
\noindent\textit{Source reference.} Original lemma 34747: sources/stacks-more-algebra/pinned-original-source.tex, source lines 34747--34747. Same-chapter navigation pointer; the statement it names is standard for the section and the argument that consumes it stays inside the counted proof span.\par\medskip
\ifdefined\hypertarget\hypertarget{sourcepointer-3}{}\fi
\noindent\textit{Source reference.} Original lemma 34641: sources/stacks-more-algebra/pinned-original-source.tex, source lines 34641--34641. Same-chapter navigation pointer; the statement it names is standard for the section and the argument that consumes it stays inside the counted proof span.\par\medskip
\ifdefined\hypertarget\hypertarget{sourcepointer-4}{}\fi
\noindent\textit{Source reference.} Original lemma 34765: sources/stacks-more-algebra/pinned-original-source.tex, source lines 34765--34765. Same-chapter navigation pointer; the statement it names is standard for the section and the argument that consumes it stays inside the counted proof span.\par\medskip
\ifdefined\hypertarget\hypertarget{sourcepointer-5}{}\fi
\noindent\textit{Source reference.} Original lemma 34490: sources/stacks-more-algebra/pinned-original-source.tex, source lines 34490--34490. Same-chapter navigation pointer; the statement it names is standard for the section and the argument that consumes it stays inside the counted proof span.\par\medskip
\begin{lemma}
\label{lemma-characteristic-p-case}
Let $A \subset B \subset C$ be extensions of discrete valuation rings
with fractions fields $K \subset L \subset M$. Assume
\begin{enumerate}
\item $A \subset B$ weakly unramified,
\item the characteristic of $K$ is $p$,
\item $M$ is a degree $p$ Galois extension of $L$, and
\item $\kappa_A = \bigcap_{n \geq 1} \kappa_B^{p^n}$.
\end{enumerate}
Then there exists a finite Galois extension $K_1/K$
totally ramified with respect to $A$
which is a weak solution for $A \to C$.
\end{lemma}

\begin{proof}
Since the characteristic of $L$ is $p$ we know that $M$ is an Artin-Schreier
extension of $L$ (Fields, Lemma \ref{fields-lemma-Artin-Schreier}).
Thus we may pick $z \in M$, $z \not \in L$ such that
$\xi = z^p - z \in L$. Choose $n \geq 0$ such that $\pi^n\xi \in B$.
We pick $z$ such that $n$ is minimal. If $n = 0$, then $M/L$ is unramified
with respect to $B$ (Lemma \ref{lemma-pre-characteristic-p-case}) and
we are done. Thus we have $n > 0$.

\medskip\noindent
Assumption (4) implies that $\kappa_A$ is perfect. Thus we may
choose compatible ring maps $\overline{\sigma} : \kappa_A \to A/\pi^n A$ and
$\overline{\sigma} : \kappa_B \to B/\pi^n B$ as in
Lemma \ref{lemma-cohen}. We lift the second of these to a
map of sets $\sigma : \kappa_B \to B$\footnote{If $B$ is complete, then
we can choose $\sigma$ to be a ring map. If $A$ is also complete and
$\sigma$ is a ring map, then $\sigma$ maps $\kappa_A$ into $A$.}.
Then we can write
$$
\xi = \sum\nolimits_{i = n, \ldots, 1} \sigma(\lambda_i) \pi^{-i} + b
$$
for some $\lambda_i \in \kappa_B$ and $b \in B$. Let
$$
I = \{i \in \{n, \ldots, 1\} \mid \lambda_i \in \kappa_A\}
$$
and
$$
J = \{j \in \{n, \ldots, 1\} \mid \lambda_i \not \in \kappa_A\}
$$
We will argue by induction on the size of the finite set $J$.

\medskip\noindent
The case $J = \emptyset$. Here for all $i \in \{n, \ldots, 1\}$ we have
$\sigma(\lambda_i) = a_i + \pi^n b_i$ for some $a_i \in A$ and $b_i \in B$
by our choice of $\sigma$. Thus
$\xi = \pi^{-n} a + b$ for some $a \in A$ and $b \in B$.
If $p | n$, then we write $a = a_0^p + \pi a_1$ for some $a_0, a_1 \in A$
(as the residue field of $A$ is perfect). We compute
$$
(z - \pi^{-n/p}a_0)^p - (z - \pi^{-n/p}a_0) =
\pi^{-(n - 1)}(a_1 + \pi^{n - 1 - n/p}a_0) + b'
$$
for some $b' \in B$. This would contradict the minimality of $n$. Thus $p$
does not divide $n$. Consider the degree $p$ extension $K_1$ of $K$ given
by $w^p - w = \pi^{-n}a$. By Lemma \ref{lemma-pre-characteristic-p-case}
this extension is Galois and totally ramified with respect to $A$.
Thus $L_1 = L \otimes_K K_1$ is a field and $A_1 \subset B_1$
is weakly unramified
(Lemma \ref{lemma-weakly-unramified-goes-up-along-totally-ramified}).
By Lemma \ref{lemma-pre-characteristic-p-case}
the ring $M_1 = M \otimes_K K_1$ is either a product of $p$ copies
of $L_1$ (in which case we are done) or a field extension of $L_1$
of degree $p$. Moreover, in the second case, either $C_1$ is weakly unramified
over $B_1$ (in which case we are done) or $M_1/L_1$ is degree $p$,
Galois, and totally ramified with respect to $B_1$.
In this last case the extension $M_1/L_1$
is generated by the element $z - w$ and
$$
(z - w)^p - (z - w) = z^p - z - (w^p - w) = b
$$
with $b \in B$ (see above). Thus by Lemma \ref{lemma-pre-characteristic-p-case}
once more the extension $M_1/L_1$ is unramified with respect to $B_1$
and we conclude that $K_1$ is a weak solution for $A \to C$.
From now on we assume $J \not = \emptyset$.

\medskip\noindent
Suppose that $j', j \in J$ such that $j' = p^r j$ for some
$r > 0$. Then we change our choice of $z$ into
$$
z' = z -
(\sigma(\lambda_j) \pi^{-j} + \sigma(\lambda_j^p) \pi^{-pj} + \ldots +
\sigma(\lambda_j^{p^{r - 1}}) \pi^{-p^{r - 1}j})
$$
Then $\xi$ changes into $\xi' = (z')^p - (z')$ as follows
$$
\xi' =
\xi - \sigma(\lambda_j) \pi^{-j} + \sigma(\lambda_j^{p^r}) \pi^{-j'}
+ \text{something in }B
$$
Writing
$\xi' = \sum\nolimits_{i = n, \ldots, 1} \sigma(\lambda'_i) \pi^{-i} + b'$
as before we find that
$\lambda'_i = \lambda_i$ for $i \not = j, j'$ and $\lambda'_j = 0$.
Thus the set $J$ has gotten smaller.
By induction on the size of $J$ we may assume no such pair $j, j'$ exists.
(Please observe that in this procedure we may get thrown back into the case
that $J = \emptyset$ we treated above.)

\medskip\noindent
For $j \in J$ write $\lambda_j = \mu_j^{p^{r_j}}$ for some $r_j \geq 0$ and
$\mu_j \in \kappa_B$ which is not a $p$th power. This is possible by our
assumption (4). Let $j \in J$ be the unique index such that $j p^{-r_j}$
is maximal. (The index is unique by the result of the preceding paragraph.)
Choose $r > \max(r_j + 1)$ and such that $j p^{r - r_j} > n$ for $j \in J$.
Choose a separable extension $K_1/K$ totally ramified with respect to $A$
of degree $p^r$ such that the corresponding discrete valuation ring
$A_1 \subset K_1$ has uniformizer $\pi'$ with
$(\pi')^{p^r} = \pi + \pi^{n + 1}a$ for some $a \in A_1$
(Lemma \ref{lemma-make-degree-q-extension}).
Observe that $L_1 = L \otimes_K K_1$ is a field and that
$L_1/L$ is totally ramified with respect to $B$
(Lemma \ref{lemma-weakly-unramified-goes-up-along-totally-ramified}).
Computing in the integral closure $B_1$ we get
$$
\xi = \sum\nolimits_{i \in I} \sigma(\lambda_i) (\pi')^{-i p^r} +
\sum\nolimits_{j \in J} \sigma(\mu_j)^{p^{r_j}} (\pi')^{-j p^r} + b_1
$$
for some $b_1 \in B_1$. Note that $\sigma(\lambda_i)$ for $i \in I$
is a $q$th power modulo $\pi^n$, i.e., modulo $(\pi')^{n p^r}$.
Hence we can rewrite the above as
$$
\xi = \sum\nolimits_{i \in I} x_i^{p^r} (\pi')^{-i p^r} +
\sum\nolimits_{j \in J} \sigma(\mu_j)^{p^{r_j}} (\pi')^{- j p^r}
+ b_1
$$
As in the previous paragraph we change our choice of $z$ into
\begin{align*}
z' & = z \\
& -
\sum\nolimits_{i \in I}
\left(x_i (\pi')^{-i} + \ldots + x_i^{p^{r - 1}} (\pi')^{-i p^{r - 1}}\right)
\\
& -
\sum\nolimits_{j \in J}
\left(
\sigma(\mu_j) (\pi')^{- j p^{r - r_j}}
+ \ldots +
\sigma(\mu_j)^{p^{r_j - 1}} (\pi')^{- j p^{r - 1}}
\right)
\end{align*}
to obtain
$$
(z')^p - z' =
\sum\nolimits_{i \in I} x_i (\pi')^{-i} +
\sum\nolimits_{j \in J} \sigma(\mu_j) (\pi')^{- j p^{r - r_j}} + b_1'
$$
for some $b'_1 \in B_1$.
Since there is a unique $j$ such that $j p^{r - r_j}$ is maximal
and since $j p^{r - r_j}$ is bigger than $i \in I$ and divisible
by $p$, we see that $M_1 / L_1$ falls into case (C) of
Lemma \ref{lemma-pre-characteristic-p-case}.
This finishes the proof.
\end{proof}

\end{document}
